\label{sec:vae-ablation}

A poorly initialised frozen encoder could severely limit the agents' language, preventing them from reaching their full cooperation potential.

\section{Discretisation Stability: VQ-VAE vs SQ-VAE}
\label{sec:vae-ablation.1}

For this evaluation, reconstruction loss and commitment loss are tracked over a 50-epoch training process using 26880 datapoints.
Liu defined a healthy encoder for the AI Mother Tongue system \cite{AIMotherTongue}.
Reconstruction loss should fall between 0.08 and 0.15.
Commitment Loss must remain below 1.0.

\begin{figure}[htbp]
    \centering
    \includegraphics[width=\textwidth]{recon_loss_eval.png}
    \caption{Reconstruction Loss over 50 epochs comparing VQ-VAE and SQ-VAE.}
    \label{fig:vae_training_eval}
\end{figure}

As shown in Figure \ref{fig:vae_training_eval}, both methods of training got below the maximum reconstruction loss of 0.15.
However, training with the SQ-VAE resulted in a final reconstruction below 0.08.
Liu's recommended loss for a healthy codebook is based on low reconstruction error, which is often a symptom of `posterior collapse' \cite{VQVAE, posterior-collapse}.
This is where the network memorises a small set of generalisable states.

To determine if this `posterior collapse' has occurred, the codebook utilisation was analysed.

\begin{table}[ht]
    \centering
    \begin{tabular}{|l|c|c|}
    \hline
    \textbf{} & \textbf{VQ-VAE} & \textbf{SQ-VAE} \\ \hline
    \textbf{Codebook Utilisation} & $66.67\% \pm 4.8\%$ & $94.8\% \pm 6.5\%$ \\ \hline
    \end{tabular}
    \caption{Codebook utilisation comparison between training with VQ-VAE and SQ-VAE.}
    \label{tab:codebook_util}
\end{table}

The table \ref{tab:codebook_util} clearly shows that using an SQ-VAE, though having a low reconstruction loss, in fact has a better codebook utilisation than using a VQ-VAE.
Therefore, this low loss indicates high model effectiveness rather than 'posterior collapse'.

By the end of training, the VQ-VAE's commitment loss had converged to $\sim 0.25$.
This is well below Liu's requirement.

\section{Codebook Initialisation}
\label{sec:vae-ablation.2}

% A potential concern of using a VQ-VAE is the previously discussed 'mode collapse'.
% The 'mode collapse' effectivly reduces the vocabulary size of the language restricting the expressiveness of the agents.
% To evaluate this concern, an ablation study is performed.
% Two initialisation strategies are compared:
% a standard uniform random initalisation,
% and a data-driven K-Means prior.

% \begin{table}[h]
%     \centering
%     \begin{tabular}{|l|c|c|c|}
%     \hline
%     \textbf{Initialisation Strategy} & \textbf{Final Validation Recon Loss} & \textbf{Active Codes} \\ \hline
%     \textbf{Random Uniform Prior} & $0.064 \pm 0.01$ & $79.13 \pm 12.63$ \\ \hline
%     \textbf{K-Means Prior} & $0.070 \pm 0.004$ & $83.53 \pm 7.67$ \\ \hline
%     \end{tabular}
%     \caption{Memory Management Overhead per Epoch: $T=20,000$, $W=32$}
%     \label{tab:init_vae_eval}
% \end{table}

% The table \ref{tab:init_vae_eval} shows the random uniform initialisation achieved a marginly lower reconstruction loss of $\sim 0.064$.
% However, the training proved to be much more volatile than the K-Means initialisation method and analysis post-training revealed a average index utilisation.
% This indicates that the VQ-VAE model initialised with random uniform embeddings is subject to the 'initialisation lottery', which can be mitigated by instead initialising codebook embeddings with a data-driven K-Means prior.
% The reduced volatility of training the language ensures high reproducibility for MARL training.

To evaluate the robustness of the VAEs, an ablation study is performed.
Two initialisation strategies are compared:
a standard uniform random initialisation,
and a data-driven K-Means prior.

For the VQ-VAE, the uniform random prior yielded superior reconstruction loss and codebook utilisation compared to the K-Means prior.
However, the SQ-VAE exhibited the opposite behaviour: the reconstruction loss and codebook utilisation were superior when using the K-Means prior rather than the random uniform prior.

\begin{figure}[H]
    \centering

    \begin{subfigure}{0.48\textwidth}
        \centering
        \includegraphics[width=\linewidth]{pca-vq-vae.png}
        \caption{VQ-VAE with prior.}
        \label{fig:pca-vq-vae}
    \end{subfigure}
    \begin{subfigure}{0.48\textwidth}
        \centering
        \includegraphics[width=\linewidth]{pca-vq-vae-no-prior.png}
        \caption{VQ-VAE without prior.}
        \label{fig:pca-vq-vae-no-prior}
    \end{subfigure}

    \begin{subfigure}{0.48\textwidth}
        \centering
        \includegraphics[width=\linewidth]{pca-sq-vae.png}
        \caption{SQ-VAE with prior.}
        \label{fig:pca-sq-vae}
    \end{subfigure}
    \begin{subfigure}{0.48\textwidth}
        \centering
        \includegraphics[width=\linewidth]{pca-sq-vae-no-prior.png}
        \caption{SQ-VAE without prior.}
        \label{fig:pca-sq-vae-no-prior}
    \end{subfigure}

    \caption{PCA projections of latent codes for SQ-VAE and VQ-VAE with and without the prior.}
    \label{fig:pca-sq-vq-vae}
\end{figure}

To visualise, Principal Component Analysis (PCA) was applied to reduce the dimensionality of the learned codebook embeddings, see Figure \ref{fig:pca-sq-vq-vae}.
This visualisation explains the findings.
For VQ-VAE (Figures \ref{fig:pca-vq-vae} and \ref{fig:pca-vq-vae-no-prior}), the random uniform initialisation scattered the codes across the latent space.
This allowed the untrained encoder to naturally interact with them, enabling gradients to flow through each embedding.
Conversely, the K-Means prior forced the embeddings into a small region, where the untrained encoder would interact only with a select few edge embeddings, leaving a high concentration of unused codes.

For (vMF) SQ-VAE (Figures \ref{fig:pca-sq-vae} and \ref{fig:pca-sq-vae-no-prior}), when using $L_2$-normalised Spherical K-Means prior, the embeddings were each set in the correct region and the model used its stochastic annealing purely for local fine-tuning.
Whereas, when using the random initialisation, embeddings frequently spawned at extreme angular distances from the true data distribution, and the model's annealing process pushed the embeddings into a suboptimal local minima.
